\textbf{Phần 1. Trắc nghiệm lựa chọn}
\setcounterpageref{ex}{0}

\begin{ex} %[4X9N50]
Biểu đồ cột thường gồm các trục nào?
\choice
{Trục số tròn và trục dọc}
{Hai trục chéo nhau}
{\True Trục nằm ngang (ghi đối tượng) và trục thẳng đứng (ghi số liệu)}
{Chỉ có một trục duy nhất}
\loigiai{
Biểu đồ cột gồm trục nằm ngang và trục thẳng đứng.\\
\textbf{Đáp án của câu hỏi là C}
}
\end{ex}

\begin{ex} %[4X9N50]
Độ cao của mỗi cột trên biểu đồ cột thể hiện
\choice
{Tên của đối tượng}
{\True Số liệu (số lượng, giá trị) của đối tượng tương ứng}
{Thời gian vẽ biểu đồ}
{Màu sắc của cột}
\loigiai{
Độ cao thể hiện số liệu của đối tượng.\\
\textbf{Đáp án của câu hỏi là B}
}
\end{ex}

\begin{ex} %[4X9H50]
Nhìn vào biểu đồ cột số cây trồng của 4 tổ: Tổ 1 ($20\text{ cây}$), Tổ 2 ($25\text{ cây}$), Tổ 3 ($15\text{ cây}$), Tổ 4 ($30\text{ cây}$). Tổ trồng nhiều cây nhất là
\choice
{Tổ 2}
{Tổ 1}
{Tổ 3}
{\True Tổ 4}
\loigiai{
Tổ 4 có cột cao nhất ứng với 30 cây.\\
\textbf{Đáp án của câu hỏi là D}
}
\end{ex}

\begin{ex} %[4X9H50]
Cũng từ biểu đồ trên, tổ trồng ít cây nhất là
\choice
{\True Tổ 3}
{Tổ 1}
{Tổ 2}
{Tổ 4}
\loigiai{
Tổ 3 có cột thấp nhất ứng với 15 cây.\\
\textbf{Đáp án của câu hỏi là A}
}
\end{ex}

\begin{ex} %[4X9H50]
Tổng số cây cả 4 tổ trồng được là
\choice
{$80\text{ cây}$}
{$100\text{ cây}$}
{\True $90\text{ cây}$}
{$85\text{ cây}$}
\loigiai{
$20 + 25 + 15 + 30 = 90$ cây.\\
\textbf{Đáp án của câu hỏi là C}
}
\end{ex}

\begin{ex} %[4X9H50]
Tổ 4 trồng nhiều hơn Tổ 3 bao nhiêu cây?
\choice
{$10\text{ cây}$}
{\True $15\text{ cây}$}
{$20\text{ cây}$}
{$5\text{ cây}$}
\loigiai{
$30 - 15 = 15$ cây.\\
\textbf{Đáp án của câu hỏi là B}
}
\end{ex}

\begin{ex} %[4X9V50]
Biểu đồ số điểm 10 trong tuần của lớp 4A: Thứ Hai ($12$), Thứ Ba ($15$), Thứ Tư ($18$), Thứ Năm ($10$), Thứ Sáu ($20$). Ngày có nhiều điểm 10 nhất là
\choice
{Thứ Tư}
{Thứ Ba}
{Thứ Hai}
{\True Thứ Sáu}
\loigiai{
Thứ Sáu có 20 điểm 10, nhiều nhất.\\
\textbf{Đáp án của câu hỏi là D}
}
\end{ex}

\begin{ex} %[4X9V50]
Tính trung bình mỗi ngày lớp 4A có bao nhiêu điểm 10 trong tuần đó?
\choice
{\True $15$ điểm 10}
{$14$ điểm 10}
{$16$ điểm 10}
{$18$ điểm 10}
\loigiai{
$(12 + 15 + 18 + 10 + 20) : 5 = 75 : 5 = 15$ điểm 10.\\
\textbf{Đáp án của câu hỏi là A}
}
\end{ex}

\begin{ex} %[4X9V50]
Biểu đồ số sách bán trong 3 tháng của một hiệu sách: Tháng 1 ($1200\text{ cuốn}$), Tháng 2 ($1500\text{ cuốn}$), Tháng 3 ($1800\text{ cuốn}$). Tổng số sách bán trong 3 tháng là
\choice
{$4200\text{ cuốn}$}
{$4800\text{ cuốn}$}
{\True $4500\text{ cuốn}$}
{$5000\text{ cuốn}$}
\loigiai{
$1200 + 1500 + 1800 = 4500$ cuốn.\\
\textbf{Đáp án của câu hỏi là C}
}
\end{ex}

\begin{ex} %[4X9V50]
Trong biểu đồ trên, số sách tháng 3 bán nhiều hơn tháng 1 là bao nhiêu cuốn?
\choice
{$500$ cuốn}
{\True $600$ cuốn}
{$700$ cuốn}
{$800$ cuốn}
\loigiai{
$1800 - 1200 = 600$ cuốn.\\
\textbf{Đáp án của câu hỏi là B}
}
\end{ex}

\begin{ex} %[4X9V50]
Biểu đồ cột biểu diễn số học sinh tham gia các câu lạc bộ: Cờ vua ($25\text{ em}$), Bơi lội ($40\text{ em}$), Bóng đá ($50\text{ em}$), Cầu lông ($35\text{ em}$). Câu lạc bộ có đông học sinh nhất là
\choice
{Bơi lội}
{Cầu lông}
{Cờ vua}
{\True Bóng đá}
\loigiai{
Bóng đá có 50 em, đông nhất.\\
\textbf{Đáp án của câu hỏi là D}
}
\end{ex}

\begin{ex} %[4X9C50]
Biểu đồ số học sinh giỏi 4 khối trường Tiểu học: Khối 1 ($80$), Khối 2 ($95$), Khối 3 ($85$), Khối 4 ($100$). Nếu muốn số học sinh giỏi trung bình mỗi khối là $95$ em thì Khối 5 cần có bao nhiêu học sinh giỏi?
\choice
{\True $115$ em}
{$110$ em}
{$120$ em}
{$105$ em}
\loigiai{
Tổng cả 5 khối cần: $95 \times 5 = 475$ em.\\
Tổng 4 khối đầu: $80 + 95 + 85 + 100 = 360$ em.\\
Khối 5 cần: $475 - 360 = 115$ em.\\
\textbf{Đáp án của câu hỏi là A}
}
\end{ex}

\textbf{Phần 2. Trắc nghiệm đúng sai}
\setcounterpageref{ex}{0}

\begin{ex} %[4X9H50]
Xét tính đúng sai khi đọc biểu đồ cột:
\choiceTF
{\True Trục nằm ngang ghi danh sách các đối tượng khảo sát.}
{\True Chiều cao của cột thể hiện số lượng của từng đối tượng.}
{Biểu đồ cột chỉ vẽ được cho tối đa 3 đối tượng.}
{\True Có thể so sánh số lượng giữa các đối tượng thông qua độ cao của các cột.}
\loigiai{
\begin{itemize}
    \item Trục ngang: Đúng.
    \item Chiều cao cột: Đúng.
    \item Biểu đồ cột có thể biểu diễn nhiều đối tượng tùy ý. Sai.
    \item So sánh qua chiều cao cột: Đúng.
\end{itemize}
\begin{itemize}
    \item a) \textbf{Đúng}.
    \item b) \textbf{Đúng}.
    \item c) \textbf{Sai}.
    \item d) \textbf{Đúng}.
\end{itemize}
}
\end{ex}

\begin{ex} %[4X9V50]
Cho biểu đồ số huy chương của 3 đội thể thao: Đội A ($12$ HCV), Đội B ($15$ HCV), Đội C ($8$ HCV). Xét tính đúng sai:
\choiceTF
{\True Đội B đạt nhiều huy chương vàng nhất.}
{\True Đội C đạt ít huy chương vàng nhất.}
{Cả 3 đội đạt tổng cộng $40$ huy chương vàng.}
{\True Cả 3 đội đạt tổng cộng $35$ huy chương vàng.}
\loigiai{
\begin{itemize}
    \item Đội B: 15 HCV (nhiều nhất). Đúng.
    \item Đội C: 8 HCV (ít nhất). Đúng.
    \item Tổng: $12 + 15 + 8 = 35$ HCV $\neq 40$ HCV. Ý 3 Sai, Ý 4 Đúng.
\end{itemize}
\begin{itemize}
    \item a) \textbf{Đúng}.
    \item b) \textbf{Đúng}.
    \item c) \textbf{Sai}.
    \item d) \textbf{Đúng}.
\end{itemize}
}
\end{ex}

\textbf{Phần 3. Trả lời ngắn}
\setcounterpageref{ex}{0}

\begin{ex} %[4X9H50]
Biểu đồ cột số bạn thích các môn thể thao: Bóng đá ($18$), Bơi ($12$), Cầu lông ($15$). Có bao nhiêu bạn thích môn bóng đá?
\shortans{18}
\loigiai{
\textbf{Đáp số:} $18$\\
18 bạn.
}
\end{ex}

\begin{ex} %[4X9H50]
Từ dữ liệu trên, tính tổng số bạn tham gia khảo sát.
\shortans{45}
\loigiai{
\textbf{Đáp số:} $45$\\
$18 + 12 + 15 = 45$ bạn.
}
\end{ex}

\begin{ex} %[4X9V50]
Biểu đồ số tấn hàng 3 xe chở: Xe 1 ($4\text{ tấn}$), Xe 2 ($6\text{ tấn}$), Xe 3 ($5\text{ tấn}$). Tính số tấn hàng trung bình mỗi xe chở.
\shortans{5}
\loigiai{
\textbf{Đáp số:} $5$\\
$(4 + 6 + 5) : 3 = 5$ tấn.
}
\end{ex}

\begin{ex} %[4X9V50]
Xe 2 chở nhiều hơn xe 1 bao nhiêu tấn hàng?
\shortans{2}
\loigiai{
\textbf{Đáp số:} $2$\\
$6 - 4 = 2$ tấn.
}
\end{ex}

\begin{ex} %[4X9V50]
Biểu đồ số bạn đọc sách: Nam ($28$), Nữ ($32$). Hỏi số bạn nữ nhiều hơn số bạn nam bao nhiêu bạn?
\shortans{4}
\loigiai{
\textbf{Đáp số:} $4$\\
$32 - 28 = 4$ bạn.
}
\end{ex}

\begin{ex} %[4X9C50]
Số học sinh 4 lớp: 4A ($32$), 4B ($35$), 4C ($33$), 4D ($36$). Tính số học sinh trung bình của mỗi lớp.
\shortans{34}
\loigiai{
\textbf{Đáp số:} $34$\\
$(32 + 35 + 33 + 36) : 4 = 136 : 4 = 34$ học sinh.
}
\end{ex}
